\documentclass[12pt,a4paper]{article}
\usepackage[a4paper,margin=18mm]{geometry}
\usepackage[T2A]{fontenc}
\usepackage[utf8]{inputenc}
\usepackage[russian]{babel}
\usepackage{amsmath,amssymb,mathtools}
\usepackage{array,booktabs,longtable}
\usepackage{xcolor}
\usepackage{hyperref}
\usepackage{tikz}

\hypersetup{
  colorlinks=true,
  linkcolor=blue!55!black,
  urlcolor=blue!55!black,
  pdftitle={Ревью домашней работы по геометрии},
  pdfauthor={Лёвин Артём Александрович}
}
\setlength{\parindent}{0pt}
\setlength{\parskip}{0.55em}
\renewcommand{\arraystretch}{1.25}

\newcommand{\tasktitle}[1]{%
  \phantomsection\label{task:#1}%
  \section*{Задание #1}%
}
\newcommand{\errorlabel}[1]{\textcolor{red}{Ошибка:} #1}
\newcommand{\keyidea}[1]{\textcolor{blue!60!black}{Ключевая идея:} #1}
\newcommand{\finalnote}[1]{\textcolor{teal!60!black}{Итог:} #1}
\newcommand{\subhead}[1]{\vspace{0.3em}\textbf{#1}\par}

\begin{document}

\begin{titlepage}
\centering
\vspace*{0.25\textheight}
{\Huge\bfseries Ревью домашней работы\par}
\vspace{1.8cm}
{\Large Автор: Лёвин Артём Александрович\par}
\vspace{0.8cm}
{\large 14 сентября 2026 года\par}
\vfill
\begin{tikzpicture}[x=1cm,y=1cm]
  \draw[blue!45!black, line width=0.9pt]
    (-5,0) .. controls (-2.4,0.55) and (2.4,-0.55) .. (5,0);
\end{tikzpicture}
\end{titlepage}

\newpage
\section*{Сводка по заданиям, требующим доработки}
\small
{\setlength{\tabcolsep}{3pt}%
\begin{longtable}{@{}>{\centering\arraybackslash}p{0.08\linewidth} p{0.18\linewidth} p{0.24\linewidth} p{0.17\linewidth} p{0.18\linewidth}@{}}
\toprule
\textbf{№} & \textbf{Тема} & \textbf{Суть ошибки} &
\textbf{Ваш ответ} & \textbf{Правильный ответ} \\
\midrule
\endfirsthead
\toprule
\textbf{№} & \textbf{Тема} & \textbf{Суть ошибки} &
\textbf{Ваш ответ} & \textbf{Правильный ответ} \\
\midrule
\endhead
\hyperref[task:3]{3} &
Равнобедренный треугольник &
Арифметическая ошибка в последнем вычислении &
$47^\circ$ &
$44^\circ$ \\

\addlinespace
\hyperref[task:5]{5} &
Первый признак равенства треугольников &
Обоснование сформулировано неточно: не перечислены все элементы, необходимые для применения признака &
$CD=7{,}4$ см &
$CD=7{,}4$ см \\

\addlinespace
\hyperref[task:6]{6} &
Медиана, высота, равнобедренный треугольник &
Решение не доведено до вычисления &
«Не помню» &
$48$ см \\

\addlinespace
\hyperref[task:9]{9} &
Доказательство свойства биссектрисы &
Верная идея сравнить треугольники, но первый признак равенства не обоснован полностью; на чертеже также без обоснования отмечен прямой угол &
$AD=DC$ &
$AD=DC$, поэтому $BD$ --- медиана \\
\bottomrule
\end{longtable}
}
\normalsize

\newpage
\thispagestyle{empty}
\vspace*{\fill}
\begin{center}
{\Huge\bfseries Разбор ошибок}
\end{center}
\vspace*{\fill}

\newpage
\tasktitle{3}

\subhead{Текст задания}
В равнобедренном треугольнике $ABC$ известно, что $AB=AC$.
Внешний угол при вершине $B$ равен $112^\circ$.
Найдите угол $A$.

\subhead{Ваше решение}
Сначала найден внутренний угол при вершине $B$:
\[
180^\circ-112^\circ=68^\circ.
\]
Затем использовано равенство углов при основании равнобедренного
треугольника, поэтому угол $C$ также принят равным $68^\circ$.
После этого записано вычисление угла $A$ через сумму углов треугольника.

\subhead{Ваш ответ}
\[
47^\circ.
\]

\subhead{Ошибка}
\errorlabel{арифметическая ошибка в последнем действии.}

\subhead{Где именно возникает ошибка}
Последний корректный шаг:
\[
\angle B=\angle C=68^\circ.
\]
Первый неверный шаг --- результат вычисления
\[
180^\circ-(68^\circ+68^\circ).
\]
Сумма $68^\circ+68^\circ$ равна $136^\circ$, поэтому из $180^\circ$
нужно вычесть $136^\circ$ и получить $44^\circ$, а не $47^\circ$.

\subhead{Почему это неверно}
Равенство $AB=AC$ означает, что основание равнобедренного треугольника
--- сторона $BC$, а углы при основании $B$ и $C$ равны.
Эта геометрическая часть выполнена верно.
Ошибка появляется только в арифметике.
Сумма внутренних углов любого треугольника равна $180^\circ$,
поэтому после нахождения двух углов третий определяется однозначно.

\subhead{Ключевая идея}
\keyidea{сначала перевести внешний угол при $B$ во внутренний,
затем использовать равенство углов при основании и только после этого
применить сумму углов треугольника.}

\subhead{Исправление от последнего правильного шага}
От шага $\angle B=\angle C=68^\circ$:
\[
\angle A=180^\circ-68^\circ-68^\circ=44^\circ.
\]

\subhead{Полное правильное решение}
Внешний угол при вершине $B$ и внутренний угол $B$ --- смежные, поэтому
\[
\angle B=180^\circ-112^\circ=68^\circ.
\]
Так как $AB=AC$, треугольник $ABC$ равнобедренный с основанием $BC$.
Следовательно,
\[
\angle B=\angle C=68^\circ.
\]
По сумме углов треугольника
\[
\angle A
=180^\circ-(68^\circ+68^\circ)
=180^\circ-136^\circ
=44^\circ.
\]

\subhead{Проверка}
\[
44^\circ+68^\circ+68^\circ=180^\circ.
\]
Кроме того, внешний угол при $B$ равен сумме двух внутренних углов,
не смежных с ним:
\[
44^\circ+68^\circ=112^\circ,
\]
что совпадает с условием.

\subhead{Итог}
\finalnote{$\angle A=44^\circ$.
Геометрический ход был выбран верно; нужно исправить только последнее
арифметическое действие.}

\subhead{Задача на закрепление}
В равнобедренном треугольнике $KLM$ известно, что $KL=KM$.
Внешний угол при вершине $L$ равен $118^\circ$.
Найдите угол $K$.

\newpage
\tasktitle{5}

\subhead{Текст задания}
Отрезки $AC$ и $BD$ пересекаются в точке $O$.
Известно, что $AO=OC$ и $BO=OD$.

а) Докажите, что треугольники $AOB$ и $COD$ равны.

б) Известно, что $AB=7{,}4$ см. Найдите $CD$.

\subhead{Ваше решение}
В записи отмечено, что равенство связано с вертикальными углами.
Далее записано
\[
\triangle AOB=\triangle COD
\quad \text{(по первому признаку)},
\]
после чего сделан вывод
\[
AB=CD=7{,}4\text{ см}.
\]

\subhead{Ваш ответ}
\[
CD=7{,}4\text{ см}.
\]

\subhead{Ошибка}
\errorlabel{обоснование равенства треугольников сформулировано недостаточно точно.}

\subhead{Где именно возникает ошибка}
Полезная и правильная идея --- использовать вертикальные углы при точке $O$.
Неточность возникает в момент применения первого признака:
для него нужно явно указать две пары равных сторон и равные углы между ними.

\subhead{Почему это неверно}
Первый признак равенства треугольников применяется, когда две стороны
и угол между ними одного треугольника соответственно равны двум сторонам
и углу между ними другого треугольника.
Здесь полный набор условий такой:
\[
AO=OC,\qquad BO=OD,\qquad \angle AOB=\angle COD.
\]
Последнее равенство верно, потому что $\angle AOB$ и $\angle COD$
--- вертикальные углы.
Именно эти три факта вместе дают первый признак равенства.

\subhead{Ключевая идея}
\keyidea{не просто назвать признак, а перечислить все соответствующие
элементы, которые позволяют его применить.}

\subhead{Исправление от последнего правильного шага}
Нужно уточнить запись о вертикальных углах:
\[
\angle AOB=\angle COD \quad \text{как вертикальные}.
\]
Вместе с условиями $AO=OC$ и $BO=OD$ получаем равенство треугольников
по двум сторонам и углу между ними.

\subhead{Полное правильное решение}
По условию
\[
AO=OC,\qquad BO=OD.
\]
При пересечении прямых $AC$ и $BD$ углы $AOB$ и $COD$ являются
вертикальными, поэтому
\[
\angle AOB=\angle COD.
\]
Следовательно,
\[
\triangle AOB\cong\triangle COD
\]
по первому признаку равенства треугольников:
по двум сторонам и углу между ними.

Из равенства треугольников следует равенство соответствующих сторон:
\[
AB=CD.
\]
Так как $AB=7{,}4$ см, то
\[
CD=7{,}4\text{ см}.
\]

\subhead{Проверка}
Стороне $AB$ первого треугольника соответствует сторона $CD$
второго треугольника:
вершине $A$ соответствует $C$, вершине $O$ --- $O$,
вершине $B$ --- $D$.
Поэтому вывод $AB=CD$ сделан по правильной паре соответствующих сторон.

\subhead{Итог}
\finalnote{числовой ответ $7{,}4$ см получен верно,
но доказательство нужно записывать полностью:
две пары равных сторон, равные вертикальные углы между ними,
затем первый признак равенства.}

\subhead{Задача на закрепление}
Отрезки $PK$ и $MN$ пересекаются в точке $S$.
Известно, что $PS=SK$, $MS=SN$ и $PM=9{,}2$ см.
Докажите равенство треугольников $PSM$ и $KSN$ и найдите $KN$.

\newpage
\tasktitle{6}

\subhead{Текст задания}
В треугольнике $ABC$ проведены медиана $BM$ и высота $BH$
к стороне $AC$.
Известно, что $AH=36$ см и $BC=BM$.
Найдите длину стороны $AC$.

\subhead{Ваше решение}
На рисунке отмечены медиана $BM$, высота $BH$ и равенство $BM=BC$,
однако дальнейший ход решения не записан.

\subhead{Ваш ответ}
«Не помню».

\subhead{Ошибка}
\errorlabel{решение не доведено до использования свойства
равнобедренного треугольника.}

\subhead{Где именно возникает ошибка}
Корректно замечены основные элементы чертежа, но после условия $BM=BC$
не сделан следующий необходимый шаг:
треугольник $BMC$ является равнобедренным.
Высота $BH$, опущенная из вершины $B$ на основание $MC$,
в таком треугольнике одновременно является медианой.

\subhead{Почему это неверно}
Если $BM=BC$, то в треугольнике $BMC$ равны две боковые стороны.
Его основание --- $MC$.
Точки $M$, $H$ и $C$ лежат на одной прямой $AC$, а $BH\perp AC$,
следовательно, $BH\perp MC$.
Значит, $BH$ --- высота равнобедренного треугольника $BMC$,
проведённая к основанию.
В равнобедренном треугольнике такая высота делит основание пополам,
поэтому
\[
MH=HC.
\]
Одновременно $BM$ --- медиана треугольника $ABC$, поэтому
\[
AM=MC.
\]
Именно эти два равенства связывают известную длину $AH$
со всей стороной $AC$.

\subhead{Ключевая идея}
\keyidea{рассмотреть треугольник $BMC$ отдельно:
из $BM=BC$ получить его равнобедренность,
а затем использовать то, что высота к основанию является медианой.}

\subhead{Исправление от последнего правильного шага}
Из $BM=BC$ следует, что $\triangle BMC$ равнобедренный.
Поэтому
\[
MH=HC.
\]
Поскольку $BM$ --- медиана треугольника $ABC$,
\[
AM=MC.
\]
Далее остаётся выразить $AH$ через части стороны $AC$.

\subhead{Полное правильное решение}
Так как
\[
BM=BC,
\]
треугольник $BMC$ равнобедренный с основанием $MC$.

Поскольку $BH$ --- высота к стороне $AC$, а $M,H,C$ лежат на $AC$,
имеем
\[
BH\perp MC.
\]
В равнобедренном треугольнике $BMC$ высота $BH$,
проведённая к основанию $MC$, является также медианой.
Поэтому
\[
MH=HC.
\]

Обозначим
\[
MH=HC=x.
\]
Тогда
\[
MC=MH+HC=2x.
\]
Так как $BM$ --- медиана треугольника $ABC$,
точка $M$ --- середина $AC$, значит
\[
AM=MC=2x.
\]
По условию $AH=36$ см, а
\[
AH=AM+MH=2x+x=3x.
\]
Следовательно,
\[
3x=36,
\]
откуда
\[
x=12.
\]
Тогда
\[
MC=24\text{ см},\qquad AM=24\text{ см},
\]
и
\[
AC=AM+MC=24+24=48\text{ см}.
\]

\subhead{Проверка}
При найденных длинах $MH=HC=12$ см, поэтому $MC=24$ см.
Так как $AM=MC$, получаем $AM=24$ см.
Тогда
\[
AH=AM+MH=24+12=36\text{ см},
\]
что совпадает с условием.

\subhead{Итог}
\finalnote{$AC=48$ см.
Ключевой переход --- увидеть равнобедренный треугольник $BMC$
и использовать свойство его высоты.}

\subhead{Задача на закрепление}
В треугольнике $ABC$ проведены медиана $BM$ и высота $BH$
к стороне $AC$.
Известно, что $AH=45$ см и $BC=BM$.
Найдите длину стороны $AC$.

\newpage
\tasktitle{9}

\subhead{Текст задания}
В равнобедренном треугольнике $ABC$ с основанием $AC$
проведена биссектриса $BD$ угла $ABC$, где $D\in AC$.
Докажите, что $BD$ является медианой треугольника $ABC$.

\subhead{Ваше решение}
Вы сравниваете треугольники $DBC$ и $ABD$ и записываете, что они равны
по первому признаку. После этого делаете вывод
\[
AD=DC.
\]
На чертеже у точки $D$ также отмечен прямой угол.

\subhead{Ваш ответ}
Сделан вывод:
\[
AD=DC,
\]
то есть получена необходимая для медианы идея.

\subhead{Ошибка}
\errorlabel{доказательство по первому признаку записано неполно,
а прямой угол при $D$ отмечен без обоснования.}

\subhead{Где именно возникает ошибка}
Идея сравнить треугольники $ABD$ и $CBD$ верная.
Первый недостаточно обоснованный шаг --- запись
\[
\triangle DBC\cong\triangle ABD
\quad \text{(по первому признаку)}
\]
без перечисления двух пар равных сторон и равных углов между ними.

Кроме того, на рисунке стоит отметка прямого угла при $D$.
Из условия сразу известно только, что $BD$ --- биссектриса и $D\in AC$;
перпендикулярность $BD$ и $AC$ в условии не дана.
Поэтому прямой угол нельзя использовать как исходный факт.

\subhead{Почему это неверно}
Первый признак равенства треугольников требует трёх конкретных фактов:
две стороны одного треугольника должны быть соответственно равны двум
сторонам другого, а углы между этими сторонами должны быть равны.

В данной задаче эти факты действительно есть:
\[
AB=BC
\]
по определению равнобедренного треугольника,
\[
BD=BD
\]
как общая сторона, и
\[
\angle ABD=\angle DBC
\]
потому что $BD$ --- биссектриса угла $ABC$.

Только после записи этих трёх равенств можно строго сослаться
на первый признак и получить равенство треугольников.

Фраза о том, что после равенства треугольников «все стороны равны»,
тоже требует уточнения: равны не все стороны между собой,
а соответствующие стороны равных треугольников.
В частности, из равенства треугольников следует именно
\[
AD=DC.
\]

\subhead{Ключевая идея}
\keyidea{для доказательства того, что $BD$ --- медиана,
достаточно строго доказать $AD=DC$.
Для этого нужно сравнить треугольники $ABD$ и $CBD$
по двум сторонам и углу между ними.}

\subhead{Исправление от последнего правильного шага}
Сохраняем правильную идею сравнить треугольники и дописываем основания:
\[
AB=BC,\qquad BD=BD,\qquad \angle ABD=\angle DBC.
\]
Следовательно,
\[
\triangle ABD\cong\triangle CBD
\]
по первому признаку равенства треугольников.
Отсюда
\[
AD=DC.
\]

\subhead{Полное правильное решение}
Рассмотрим треугольники $ABD$ и $CBD$.

Так как треугольник $ABC$ равнобедренный с основанием $AC$, то
\[
AB=BC.
\]
Сторона $BD$ является общей:
\[
BD=BD.
\]
По условию $BD$ --- биссектриса угла $ABC$, поэтому
\[
\angle ABD=\angle DBC.
\]
Следовательно,
\[
\triangle ABD\cong\triangle CBD
\]
по первому признаку равенства треугольников:
по двум сторонам и углу между ними.

Из равенства треугольников следует равенство соответствующих сторон:
\[
AD=DC.
\]
Значит, точка $D$ является серединой стороны $AC$.
Отрезок $BD$ соединяет вершину $B$ с серединой противоположной стороны
$AC$, следовательно, по определению
\[
BD
\]
является медианой треугольника $ABC$.

\subhead{Проверка}
Для вывода о медиане требуется ровно то, что точка $D$ является серединой
стороны $AC$.
Это доказано равенством
\[
AD=DC.
\]
Никакое предположение о прямом угле при $D$ для доказательства не требуется.

\subhead{Итог}
\finalnote{основная идея решения верная и вывод $AD=DC$ выбран правильно.
Нужно полностью обосновать первый признак:
$AB=BC$, $BD$ --- общая сторона,
$\angle ABD=\angle DBC$ как половины угла при биссектрисе.
Отметку прямого угла при $D$ без отдельного обоснования ставить не нужно.}

\subhead{Задача на закрепление}
В равнобедренном треугольнике $KMN$ с основанием $KN$
проведена биссектриса $MP$ угла $KMN$, где $P\in KN$.
Докажите, что $MP$ является медианой треугольника $KMN$.

\vfill
\begin{center}
\small Лёвин Артём Александрович эксклюзивно для Екатерины
\end{center}

\end{document}
